% HAND-WRITTEN fragment. Arithmetic behind every number here, so it can be
% re-derived:
%   quantisation sigma  = 1.591 mV / sqrt(12) = 0.459 mV; difference of two
%                         independent conversions = 0.650 mV
%   potential residual  = RSS(0.650 offset, 0.931 droop fit, 0.788 ladder INL)
%                       = 1.38 mV rms; with a per-code table 1.31 mV
%   current residual    = RSS(0.650 LSB zero, 0.294 LSB converter residual)
%                       = 0.71 LSB
%   range condition     = 0.815 V / (3 x 5.7176 mV) = 47.5
%   uncalibrated yields recomputed from the published mean/sigma against
%   4 LSB = 6.364 mV (see note_calibration_intro).

\subsubsection{Residual, and what calibration cannot reach}
\label{sss:calibration-residual}

Each zero is the difference of two conversions, so its floor is two
quantisations of the converter's \SI{1.591}{\milli\volt} step:
\SI{0.65}{\milli\volt} rms. Averaging reaches that floor only if the input
dithers. It does not on its own --- with the electrode isolated the amplifier's
output noise is tens of \si{\micro\volt}, a fiftieth of a code, so repeated
conversions return the same number --- so the common-mode DAC supplies the
dither: three of its \SI{529}{\micro\volt} steps span one converter code.

On the potential axis the residual is that offset floor, the reference-droop
fit residual of \SI{0.93}{\milli\volt} rms, and the ladder's own INL of
\SI{0.79}{\milli\volt} left uncorrected, which combine to
\SI{1.38}{\milli\volt} rms against \SI{5.71}{\milli\volt} of uncalibrated
\(\sigma\). Per-channel yield inside \(\pm4\,\mathrm{LSB}\) goes from
\SI{73.2}{\percent} to effectively unity, and four-channel yield from
\SI{28.6}{\percent}. On the current axis the zero residual is
\SI{0.65}{LSB}, and with the converter's \SI{0.29}{LSB} rms departure from its
own fitted line the total is \SI{0.71}{LSB} plus \SIrange{0.1}{0.4}{\percent} of
reading in gain --- against \SIrange{12.1}{167.3}{LSB} of uncalibrated
\(\sigma\). The gain term is the tester's accuracy plus the ladder matching
wherever a gain code was extrapolated rather than measured.

\paragraph{The condition under which the current zero is correctable at all.}
A digital constant removes an offset only while the offset stays inside the
converter's window. The zero-current offset is a \emph{current}, the loop offset
divided by the electrode impedance, so what it costs at the converter scales
with \(R_f\): the requirement is \(3\sigma(v_{\mathrm{os}}) \cdot R_f /
R_{\mathrm{cell}} < \SI{0.815}{\volt}\), or
\[
  \frac{R_f}{R_{\mathrm{cell}}} \;<\; 47.5 .
\]
The Monte Carlo of Section~\ref{sss:pixel-mc} sits at 56 --- the
\SI{560}{\kilo\ohm} tap against a \SI{10}{\kilo\ohm} electrode --- which is why
some of its samples rail, and those samples are not recoverable by any constant.
Every operational pairing in this chapter is far inside the limit: the
\SI{1}{\mega\ohm} electrode at the top gain code is a ratio of 1.9, a margin of
\num{24}. The rule belongs in the gain-code selection, not in the calibration.

\paragraph{What no per-chip constant reaches.} The reference DAC's
\(-6.6\,\mathrm{LSB}\) of DNL is lost command resolution, up to
\SI{3.5}{\milli\volt} at the worst code, and a table can only choose the best
available code, not create a missing one. The converter's linearity is not
measured --- the sweep behind Section~\ref{sss:adc-transfer} steps 39 codes at a
time --- and the chip carries no ramp source precise enough to measure it, so it
is a bound and not a correction. The converter's lost \SI{34.9}{\percent} of
input range is a constraint on \(R_f \cdot I_{\mathrm{FS}}\), not an error.
Bottom-code drive nonlinearity, 1.26 LSB at \(N = 0\) and \(\leq\)\,0.004 LSB
everywhere else, varies with the current being measured and a two-point gain
constant cannot see it.

Electrode-to-electrode variation is not electronics and none of it appears in
any Monte Carlo run in this chapter, because the electrode is a fixed model.
Three of its terms behave differently under calibration. A reference-electrode
junction potential is indistinguishable from A1's offset, so a zero taken
\emph{wet} against that electrode absorbs both --- and a zero taken dry, as
every step above does, absorbs neither. Electrode area and kinetics set the
current-per-concentration slope, which only a solution of known concentration
can fix; nothing on this chip can. Uncompensated electrolyte resistance appears
as \(i\!\cdot\!R_u\) on the potential axis and \emph{is} correctable, on the same
arithmetic as the working-electrode switch drop, provided an impedance
measurement has supplied \(R_u\) for that site.

\paragraph{Open.} The temperature dependence of the constants is unquantified:
no block in this chapter has a temperature-resolved Monte Carlo, so how far a
power-up re-measurement tracks a \SI{37}{\celsius} operating point is not known.
Neither is the chip-to-chip spread of the converter's gain, which rests on one
nominal extracted run (Section~\ref{ss:adc}) and enters every current reading
directly. And none of this has been simulated: the channel has never been
netlisted with the converter in its loop, so the end-to-end sweep that would
verify the residuals above does not exist.
